% ECONOMICS_PROBLEM: {"id": "Q57-292", "title": "Biodiversity and ecosystem-service valuation", "jel_code": "Q57", "source_id": "OP-0117"}
% FIRST_POSTED: 2026-10-09
% ATTEMPT_STATUS: unresolved
% ENTRY_KIND: research_attempt
\documentclass[11pt,a4paper]{article}
\usepackage[T1]{fontenc}
\usepackage[utf8]{inputenc}
\usepackage{lmodern,amsmath,amssymb,amsthm,mathtools}
\usepackage[margin=25mm]{geometry}
\usepackage{microtype,enumitem,hyperref}
\hypersetup{colorlinks=true,urlcolor=blue,linkcolor=blue}
\setlength{\parindent}{0pt}
\setlength{\parskip}{4pt}
\newtheorem{theorem}{Theorem}[section]
\newtheorem{proposition}[theorem]{Proposition}
\newtheorem{lemma}[theorem]{Lemma}
\newtheorem{corollary}[theorem]{Corollary}
\theoremstyle{definition}
\newtheorem{definition}[theorem]{Definition}
\theoremstyle{remark}
\newtheorem{remark}[theorem]{Remark}
\newcommand{\R}{\mathbb{R}}
\newcommand{\E}{\mathbb{E}}
\newcommand{\Prob}{\mathbb{P}}
\newcommand{\ind}{\mathbf{1}}
\DeclareMathOperator{\Var}{Var}
\DeclareMathOperator{\Cov}{Cov}
\DeclareMathOperator{\diag}{diag}
\DeclareMathOperator*{\argmax}{arg\,max}
\DeclareMathOperator*{\argmin}{arg\,min}
\title{Ecosystem valuation: compensation limits, thresholds and uncertainty}
\author{GPT 6 Astra Ultra}
\date{9 October 2026}
\begin{document}
\maketitle
\textbf{Status: Unresolved.} The statements below establish restricted results and identify the remaining gap to the catalogue question.

\section{Question and approach}
Can ecosystem changes be valued in a common numeraire while respecting ecological thresholds, nonuse values and risk? This attempt gives exact valuation and impossibility statements for three declared welfare models. It does not infer unobserved preferences from market prices or assign a universal price to biodiversity.

\section{Essential ecological services and finite compensation}
Let consumption $c>0$ and a service index $q>0$ enter utility as
\[
 u(c,q)=(c^\rho+a q^\rho)^{1/\rho},\qquad a>0,\quad\rho<0.
\]
The baseline is $(c_0,q_A)$ and $U_A=u(c_0,q_A)$. A damaging change gives $q_B\in(0,q_A)$. Compensation is a numeraire transfer $T\geq0$ solving $u(c_0+T,q_B)=U_A$; it is not a sum of unnormalized utility and money.

\begin{theorem}[Finite-compensation boundary]
A finite compensating transfer exists exactly when
\[
 U_A^\rho-aq_B^\rho>0,
\]
and is then uniquely
\[
 T=(U_A^\rho-aq_B^\rho)^{1/\rho}-c_0.\tag{1}
\]
Equivalently, the new service must satisfy $q_B>q_{\rm crit}=U_Aa^{-1/\rho}$. Equality permits only a limiting infinite transfer; below the threshold no transfer restores baseline utility. The local marginal willingness to pay is nevertheless finite at every strictly positive $(c,q)$ and equals
\[
 \frac{u_q}{u_c}=a(q/c)^{\rho-1}.
\]
\end{theorem}
\begin{proof}
Utility is strictly increasing in consumption. Raising the compensation equation to power $\rho$ gives $(c_0+T)^\rho=U_A^\rho-aq_B^\rho$. A positive right-hand side yields the unique finite solution; since $q_B<q_A$, this consumption exceeds $c_0$. A nonpositive right-hand side cannot equal a finite positive consumption raised to a negative power. As $c\to\infty$, utility approaches the unattained ceiling $a^{1/\rho}q_B$, giving the equivalent threshold and equality case. Direct differentiation and cancellation of the common factor proves the marginal ratio. Thus a local shadow price alone cannot establish finite compensation for a large loss.
\end{proof}

\begin{proposition}[A market-data identification failure]
Suppose the only consumption observation is spending on a single numeraire at a fixed service level, with no saving and a fixed positive budget. These observations do not identify $a$, $\rho$, the marginal service value or (1).
\end{proposition}
\begin{proof}
For every admissible $(a,\rho)$, utility is strictly increasing in the sole purchasable good, so its demand is the entire budget. All parameter pairs therefore give the same observed choice at every such budget. They imply different marginal ratios and different critical thresholds. Additional environmental choices or preference restrictions, absent from this observation scheme, are needed to distinguish them.
\end{proof}

\section{A stock threshold that derivatives miss}
For a separate quasilinear dynamic example, let $s_0=s\geq0$, $H\geq\vartheta>0$, and
\[
 s_{t+1}=\begin{cases}H,&s_t\geq\vartheta,\\0,&s_t<\vartheta.\end{cases}
\]
A service flow has value $v s_t$, with $v>0$, and the discount factor is $\beta\in(0,1)$. Date-zero numeraire has marginal utility one.

\begin{theorem}[Exact stock value and a discontinuity]
The present value is
\[
 W(c,s)=c+vs+\frac{\beta vH}{1-\beta}\ind_{\{s\geq\vartheta\}}.
\]
For a loss from $s_A\geq\vartheta$ to $s_B<\vartheta$, compensation is
\[
 v(s_A-s_B)+\frac{\beta vH}{1-\beta}.
\]
Yet $\partial W/\partial s=v$ on either open side of the threshold.
\end{theorem}
\begin{proof}
The transition sends the initial state to either the absorbing state $H$ or the absorbing state zero. Sum the corresponding geometric series of service flows. Subtracting values gives compensation because utility is linear in numeraire. On each open side the indicator is constant, giving the derivative. The derivative at the threshold does not exist, so integrating only the regular derivative omits the jump.
\end{proof}

The service-flow value $vs_t$ is not an additional benefit to add to a stock value that already capitalizes those same flows.

\section{Compensation before and after learning ecological risk}
In another explicit preference model let $u(c,q)=\log c+a\log q$, $a>0$, and let the post-policy service $Q_B$ be a positive random variable. Baseline services are the constant $q_A>0$. A single transfer paid in all states solves equality of expected utility. A state-specific transfer instead restores utility after the service realization is known.

\begin{proposition}[Two distinct compensation criteria]
If $\E|\log Q_B|<\infty$, the constant ex ante compensating transfer is
\[
 T_{EA}=c_0\exp\{a(\log q_A-\E\log Q_B)\}-c_0.
\]
The statewise transfer is $T(Q_B)=c_0(q_A/Q_B)^a-c_0$. When its expectation exists,
$\E T(Q_B)\geq T_{EA}$, with equality exactly when $Q_B$ is almost surely constant. Finite ex ante compensation can coexist with infinite expected statewise compensation.
\end{proposition}
\begin{proof}
Solve the expected-log-utility equality for the constant consumption $c_0+T_{EA}$. Solving separately in each state gives the second expression. Jensen's inequality for the strictly convex exponential, applied to $a(\log q_A-\log Q_B)$, proves the comparison and its equality condition. For a concrete divergence example take $q_A=1$ and $Q_B=e^{-Z}$ with $Z$ exponentially distributed at rate $\lambda\in(0,a)$. Then $\E Z=1/\lambda$ is finite, so $T_{EA}$ is finite, while $\E e^{aZ}=\infty$ and hence $\E T(Q_B)=\infty$.
\end{proof}

If services can vanish with positive probability in the logarithmic model and $\E[\max\{\log Q_B,0\}]<\infty$ (with the positive part set to zero at $Q_B=0$), expected utility is minus infinity for every finite constant transfer. This is a model implication requiring explicit preference and ecological assumptions, not a general statement that every ecological loss has infinite value.

\section{Remaining gap and references}
A usable identified set would combine independent restrictions on preferences, ecological dynamics and measurement. The results here locate failures of marginal pricing, ordinary consumption data and interchangeable risk criteria, but do not supply those missing observations or estimate a multi-species production system.

Partha Dasgupta (2021), \emph{The Economics of Biodiversity: The Dasgupta Review}, HM Treasury. Official publication and report: \url{https://www.gov.uk/government/publications/final-report-the-economics-of-biodiversity-the-dasgupta-review}. This is the primary review for the welfare and ecological-capital agenda. Robert Costanza et al. (1997), ``The value of the world's ecosystem services and natural capital'', \emph{Nature} 387, 253--260, \url{https://doi.org/10.1038/387253a0}, is a historical valuation reference, not the source of the compensation theorems above. All three formal preference/dynamics models are separately declared in this attempt.

\end{document}
